\section{Evaluación y resultados}

\subsection{Por qué se construyó un banco de pruebas}

Un asistente conversacional puede dar una respuesta equivocada con total
naturalidad. A diferencia de un error de programación, que se manifiesta con
una excepción, una cifra incorrecta expresada en una frase bien construida
resulta indistinguible de una correcta para quien la escucha.

Durante el desarrollo se detectó un caso que ilustra el problema. Ante la
pregunta ``¿he gastado más este año que el año pasado?'', el modelo generaba una
consulta con \texttt{strftime('\%Y-\%m')} en lugar de \texttt{strftime('\%Y')}.
La consulta comparaba el mes actual contra el mes anterior, pero la respuesta
los presentaba como años. El asistente respondía que se había gastado más este
año cuando la realidad era 17.870 € frente a 26.643 €: la conclusión estaba
invertida y se expresaba con seguridad. Revisando la interfaz, la respuesta
resultaba convincente.

Por este motivo se desarrolló \texttt{tests/evaluar.py}, un banco de preguntas
que se ejecuta contra el agente real y mide su comportamiento de forma
sistemática.

\subsection{Qué se mide y por qué se separa}

El banco contiene 44 preguntas etiquetadas, distribuidas entre los tipos de
consulta que menciona el enunciado (gastos, comparativas y suscripciones) y
otros añadidos durante el desarrollo: seguimiento multiturno, operaciones de
Bizum, saldo, ingresos, preguntas ambiguas, límites del asistente y previsión.

Se puntúan cuatro indicadores por separado, porque responden a fallos de
naturaleza distinta:

\begin{itemize}
    \item \textbf{Elección de herramienta.} Preguntar por las suscripciones y
    que el modelo escriba una consulta SQL es un fallo aunque la cifra final sea
    correcta. Se registra mediante un espía sobre \texttt{ejecutar\_tool}, de
    forma que también se observan las resoluciones directas del backend, que no
    aparecen en el historial del modelo.

    \item \textbf{Respuesta final.} Que la consulta sea correcta no garantiza
    que el modelo traslade bien la cifra al texto. Es lo único que oye el
    usuario, y se comprueba que los valores esperados aparezcan con formato
    español.

    \item \textbf{Refuerzo visual.} Se puntúa en los dos sentidos: no mostrar
    nada cuando hay varios valores comparables es un fallo, y mostrar algo
    cuando la respuesta es un único dato también. Una tabla cuenta igual que un
    gráfico, porque la elección entre ambos corresponde al modelo y las dos son
    respuestas válidas.

    \item \textbf{Validez de la especificación.} Que llegue un gráfico no
    significa que se vea. Una especificación sin el campo \texttt{mark} supera
    la validación del backend y falla después en el renderizado, sin dejar
    rastro. Se puntúa por separado del indicador anterior porque son fallos de
    cosas distintas: uno mide al modelo decidiendo y el otro al modelo
    redactando.
\end{itemize}

Se mide además un quinto indicador sin puntuación, como diagnóstico: se ejecutan
la consulta generada por el agente y una de referencia escrita a mano, y se
comparan los valores obtenidos. No cuenta para el resultado porque el modelo
puede elegir una interpretación distinta y defendible —otro periodo, otra forma
de calcular una media— y divergir de la referencia sin estar equivocado. Nunca
se compara el texto de la consulta: existen muchas consultas correctas
distintas para la misma pregunta.

La latencia se mide en dos tramos. El relevante es el tiempo hasta el evento
\texttt{fin\_respuesta}, que es el que activa la síntesis de voz y por tanto lo
que el usuario percibe como demora. El motor visual se ejecuta después de ese
evento de forma deliberada, así que contabilizar únicamente el tiempo total
penalizaría un trabajo que ya no bloquea la respuesta.

\subsection{Resultados}

Las mediciones se realizan con tres repeticiones de las 44 preguntas, sobre
\texttt{qwen3:8b} y con la base de datos regenerada previamente.

\begin{table}[h]
\centering
\begin{tabular}{lr}
\toprule
\textbf{Indicador} & \textbf{Resultado} \\
\midrule
Elección de herramienta            & 97,6 \% \\
Respuesta final correcta           & 88,9 \% \\
Refuerzo visual cuando corresponde & 100 \% \\
Especificaciones válidas           & 100 \% \\
Latencia mediana hasta la voz      & 3,3 s \\
\bottomrule
\end{tabular}
\caption{Resultados sobre 44 preguntas y 3 repeticiones.}
\end{table}

La variabilidad entre ejecuciones del mismo sistema es de uno o dos casos sobre
126, por lo que las comparaciones entre configuraciones se realizan siempre con
repeticiones. Una diferencia de dos casos en una única ejecución no constituye
una señal.

\subsection{Decisiones de diseño tomadas a partir de la medición}

El banco de pruebas no se utilizó únicamente para verificar el resultado final.
En varios casos, la medición contradijo una decisión previa y motivó un cambio
de arquitectura.

\subsubsection{La decisión de mostrar un gráfico}

Inicialmente, la decisión de acompañar una respuesta con un gráfico se dejaba al
modelo mediante instrucciones en el \textit{system prompt}. La medición mostró
que esa decisión no era estable: variaba ante cualquier modificación del
\textit{prompt}, incluso ante cambios que no trataban sobre gráficos.

\begin{table}[h]
\centering
\begin{tabular}{lcccc}
\toprule
\textbf{Pregunta} & \textbf{v1} & \textbf{v2} & \textbf{v3} & \textbf{v4} \\
\midrule
Ranking de comercios          & 0/3 & 2/3 & 0/3 & 0/3 \\
Comparación mensual           & 0/3 & 0/3 & 0/3 & 0/3 \\
Comparación por categorías    & 0/3 & 2/3 & 2/3 & 0/3 \\
Serie de seis meses           & 0/3 & 3/3 & 3/3 & 0/3 \\
Comparación interanual        & 0/3 & 0/3 & 0/3 & 0/3 \\
Evolución mensual             & 3/3 & 3/3 & 3/3 & 3/3 \\
\bottomrule
\end{tabular}
\caption{Veces que cada pregunta produce un gráfico, sobre cuatro versiones del
\textit{system prompt}. Entre la tercera y la cuarta solo cambiaron dos líneas
relativas al tratamiento de fechas.}
\end{table}

Dentro de cada versión el comportamiento es prácticamente determinista, pero
entre versiones cambia sin relación con el contenido modificado. Cinco de las
seis preguntas no producían gráfico en la versión final; la única que funciona
de forma consistente es la que contiene una petición visual explícita
(``muéstrame la evolución'').

A partir de este resultado se reorganizó la responsabilidad. La decisión de
\textit{si} existe algo que representar pasó al backend, mediante una regla
determinista sobre la forma del resultado: varias filas con alguna columna
numérica, o una sola fila con dos o más cifras comparables. La decisión de
\textit{qué} representación utilizar —tabla o gráfico—, con qué marca, qué ejes
y con qué justificación, permanece en el modelo, en una llamada dedicada de
tarea única.

El indicador pasó del 16,7 \% al 100 \%, sin variación en la precisión de las
respuestas ni en la latencia percibida. Conviene subrayar que la decisión que se
trasladó al código no es una decisión de diseño visual, sino una comprobación
mecánica sobre la estructura de los datos.

\subsubsection{El tamaño del \textit{system prompt}}

Durante el desarrollo, el \textit{system prompt} fue creciendo al incorporar
nuevas reglas. La medición mostró que ese crecimiento degradaba reglas
anteriores: una pregunta de comparación entre categorías pasó de 3/3 a 0/3 a lo
largo de varias sesiones, con el ejemplo que la resolvía intacto en el
\textit{prompt}.

La causa era la duplicación: las mismas seis herramientas se describían tanto en
el \textit{prompt} como en sus esquemas JSON, y ambos textos se envían en cada
petición. Se eliminó la parte duplicada —la que indica cuándo llamar a cada
herramienta, que el esquema ya expresa y en mejor posición, junto a la
definición de la función— y se conservó la que el esquema no puede expresar:
cómo interpretar el resultado devuelto.

\begin{table}[h]
\centering
\begin{tabular}{lrr}
\toprule
 & \textbf{Antes} & \textbf{Después} \\
\midrule
Contexto fijo por petición   & 3.822 tokens & 3.208 tokens \\
Porcentaje de la ventana     & 47 \%        & 39 \% \\
Elección de herramienta      & 97,6 \%      & 97,6 \% \\
Respuesta final correcta     & 85,7 \%      & 90,5 \% \\
\bottomrule
\end{tabular}
\caption{Efecto de eliminar las instrucciones duplicadas.}
\end{table}

Se verificó una por una que ninguna de las diecisiete reglas obtenidas mediante
medición se hubiera perdido en el proceso. Se recuperaron dos preguntas, una de
ellas registrada desde el inicio como limitación asumida del modelo: no lo era.

\subsection{Validación de los componentes deterministas}

Los componentes que no dependen del modelo se validan por separado, con pruebas
que no requieren el LLM y que por tanto son deterministas y rápidas.

\subsubsection{Detección de pagos recurrentes}

Los umbrales del detector se ajustaron observando una única base de datos, lo
cual no demuestra que generalicen. Para comprobarlo se regenera el histórico con
500 semillas distintas y se contrasta el resultado con la verdad conocida por
construcción: 0 falsos negativos, 1 falso positivo (0,2 \% de los históricos
generados), ninguna cadencia mal clasificada y el seguro anual detectado en las
500 semillas.

Dos decisiones de diseño provienen de esa medición. La primera, reducir el
umbral de irregularidad de 0,10 a 0,06. La segunda, ampliar el histórico
generado de 15 a 24 meses, que por sí sola redujo los falsos positivos de 23 a
0: con más histórico, los comercios no recurrentes acumulan cargos suficientes
para que su irregularidad resulte evidente.

Conviene señalar el límite de esta validación. El generador asigna a los pagos
recurrentes un día fijo de cada mes, mientras que los gastos no recurrentes caen
en un día aleatorio. La prueba demuestra por tanto que el detector distingue un
cargo de día fijo de uno de día aleatorio, y que no clasifica como recurrente un
gasto habitual: una compra semanal, incluso perfectamente regular, no se detecta
como suscripción, porque el clasificador solo admite cadencias reconocidas. No
cubre, en cambio, el desplazamiento del día de cobro habitual en las
domiciliaciones reales, que se cargan el siguiente día hábil. Medido
posteriormente, el detector reconoce la cadencia con desplazamientos de hasta
dos días, y deja de reconocerla a partir de tres.

\subsubsection{Proyección de gasto}

El estimador de la proyección a fin de mes se seleccionó comparando tres
alternativas sobre los 23 meses completos del histórico, midiendo el error de
cada una según el día del mes en que se realiza la estimación.

\begin{table}[h]
\centering
\begin{tabular}{lrrrr}
\toprule
\textbf{Método} & \textbf{Día 5} & \textbf{Día 10} & \textbf{Día 15} & \textbf{Día 20} \\
\midrule
Proporción directa            & 174,5 \% & 73,3 \% & 46,1 \% & 25,9 \% \\
Pagos fijos por separado      & 31,7 \%  & 18,4 \% & 13,6 \% & 10,6 \% \\
Combinado con media histórica & 8,7 \%   & 7,6 \%  & 8,4 \%  & 7,7 \% \\
\bottomrule
\end{tabular}
\caption{Error medio de cada estimador según el día del mes.}
\end{table}

La proporción directa resulta inservible porque el alquiler y los recibos
domiciliados se cargan a principios de mes: el día 3, con 890,34 € gastados de
los cuales 844,80 € corresponden a pagos fijos, una regla de tres proyectaría
8.903 € para el mes completo.

El estimador adoptado no extrapola los pagos fijos, cuyo coste mensual ya conoce
el detector de recurrencia, y combina el ritmo de gasto variable observado con
la media histórica, ponderando según lo avanzado que esté el mes. El error se
mantiene estable en torno al 8 \% durante todo el mes, lo que permite ofrecer
una previsión útil desde los primeros días.

Cada proyección se acompaña de su margen de error, calculado aplicando el mismo
estimador a cada mes pasado y midiendo la desviación. El margen varía
considerablemente según la categoría: 0 \% en alquiler o gimnasio, que son pagos
fijos conocidos; en torno al 20 \% en supermercado o gasolina; y superior al
100 \% en los Bizums enviados, que son impredecibles por naturaleza. Ofrecer una
cifra única para todas las categorías habría resultado igual de creíble y
sustancialmente menos honesto.

\subsubsection{Flujo de operaciones}

El envío de Bizums se valida mediante 35 comprobaciones que recorren el flujo
completo sin intervención del modelo: reconocimiento de la petición, validación
del importe, resolución del contacto, gestión del PIN y atomicidad de la
escritura.

Esta prueba se escribió antes de refactorizar el flujo, precisamente porque el
banco de preguntas no cubre ese camino: sus preguntas sobre Bizums son consultas
sobre envíos pasados, que se resuelven mediante SQL. Era posible alterar el
comportamiento del envío sin que ninguna métrica lo reflejara.

Una de las comprobaciones verifica que el bloqueo de escritura se adquiere antes
de la primera lectura. Sin esa garantía, dos envíos simultáneos leen el mismo
saldo, ambos superan la comprobación y una actualización sobrescribe a la otra:
en la reproducción del caso, dos envíos de 600 € sobre un saldo de 1.000 €
resultaban ambos aprobados. La comprobación se formula sobre la propiedad que lo
impide y no sobre el síntoma, ya que una condición de carrera solo se manifiesta
cuando dos hilos coinciden en una ventana de milisegundos y una prueba que
dependiera de ello resultaría inestable.

\subsection{Limitaciones conocidas}

La evaluación permite delimitar con precisión lo que el sistema no resuelve:

\begin{itemize}
    \item \textbf{Resolución de referencias entre turnos.} Ante la secuencia
    ``¿cuánto he gastado en supermercado este año?'', ``¿y en restaurantes?'',
    ``¿y el año pasado?'', el modelo resuelve la tercera pregunta contra la
    primera y responde sobre supermercado. Es un comportamiento sistemático, no
    ocasional.

    \item \textbf{Previsión de meses futuros.} La herramienta proyecta
    únicamente el mes en curso. Ante una pregunta sobre el mes siguiente, el
    modelo responde con la cifra del mes actual sin advertir la diferencia. Se
    probaron dos correcciones —una regla en el \textit{prompt} y un campo
    explícito en la respuesta de la herramienta— y ninguna resultó efectiva.

    \item \textbf{Preguntas sin periodo explícito y consultas sobre terceros.}
    Se documentan como limitaciones asumidas, con la salvedad de que una de las
    inicialmente clasificadas como tales resultó no serlo: desapareció al
    reducir el tamaño del \textit{system prompt}.
\end{itemize}

Delimitar estas limitaciones con precisión resulta más útil que un resultado
global elevado, ya que indica en qué condiciones conviene confiar en el sistema
y en cuáles no.
